% ------------------------------------------------------------
% CHAPTER 12
% ------------------------------------------------------------

\chapter{Delta--Sigma Modulation as Measurement and Reconstruction}

The delta--sigma modulator has traditionally been understood as a circuit or algorithm that converts an analog or digital signal into a high-frequency bitstream whose long-term average reconstructs the input. While this description is adequate for engineering design, it conceals the deeper mathematical structure of the process. At its core, delta--sigma modulation is a measurement--reconstruction cycle. The quantizer acts as a measurement device that performs a symbolic collapse of the state, while the loop filter reconstructs a smooth estimate from this collapsed representation. The modulator therefore alternates between two fundamental operations: measurement, which injects entropy through symbolic collapse, and reconstruction, which reduces entropy by smoothing the internal state. This dual structure mirrors the classical description of a measurement process in quantum mechanics \cite{vonneumann1932foundations}, as well as the operational and generalized-measurement perspective developed by Busch, Lahti, and Mittelstaedt \cite{busch1996quantummeasurement}. In this chapter we formalize this analogy by interpreting quantization as a deterministic classical measurement operator and the loop filter as a reconstruction map that plays the role of a conditional expectation.

To establish the measurement-theoretic foundations, let the state of the modulator at time $n$ be $x[n] \in \mathbb{R}^{L}$, and let the quantizer be a map
\[
Q : \mathbb{R}^{L} \to \Lambda,
\]
where $\Lambda = \{q_{0}, q_{1}, \ldots, q_{M-1}\}$ is a finite set of discrete symbols. The set $\mathbb{R}^{L}$ is partitioned into $M$ regions
\[
\Omega_{k} = Q^{-1}(q_{k}),
\]
so that $Q$ is constant on each region. Each region $\Omega_{k}$ is measurable and the collection $\{\Omega_{k}\}$ forms a measurable partition of $\mathbb{R}^{L}$. This structure is equivalent to a projection-valued measure in classical measurement theory. Unlike quantum measurements, which are stochastic even for pure states, the delta--sigma quantizer is deterministic: the symbol is uniquely determined by the state. Nevertheless, the structure of the map parallels the collapse associated with a von Neumann measurement, where the continuous state $x[n]$ is compressed into a symbolic outcome $q_{k}$.

The modulator completes the measurement--reconstruction cycle when the loop filter applies the update
\[
x[n+1] = F(x[n], u[n], Q(x[n])),
\]
where $F$ is a vector field that smooths the effects of collapse generated by $Q$. This vector field plays the role of an estimator, updating the internal state to restore coherence after symbolic collapse. This reconstruction step is analogous to Bayesian or variational state estimation, as the modulator reconstructs an internal estimate that remains consistent with the measurement outcome and the input. The alternation between collapse and reconstruction yields the characteristic entropy-shaping behavior of the modulator.

To formalize this structure mathematically, we begin with a precise definition of measurement.

\begin{definition}
A \emph{deterministic measurement} on a measurable space $(X,\Sigma)$ is a measurable map $M : X \to \Lambda$ where $\Lambda$ is a finite set equipped with the discrete sigma-algebra. The associated \emph{instrument} is the map $\mathcal{I} : \Sigma \times \Sigma_{\Lambda} \to [0,1]$ defined by
\[
\mathcal{I}(A \mid \lambda) = \mu(A \cap M^{-1}(\lambda)),
\]
where $\mu$ is a measure on $X$.
\end{definition}

This notion generalizes the classical part of the operational approach to measurement \cite{busch1996quantummeasurement}. In the case of a delta--sigma modulator, $X = \mathbb{R}^{L}$, $\Sigma$ is the Borel sigma-algebra, and $M = Q$ is exactly the quantizer. The quantizer thus defines a deterministic classical measurement in the sense of generalized measurement theory.

The collapse of the state under measurement is represented by the map
\[
x[n] \longmapsto Q(x[n]),
\]
which forgets all information within each region $\Omega_{k}$ except for the label $q_{k}$. In quantum theory this collapse is expressed by a projection operator, whereas in the classical case the collapse is effected by taking the coset of the partition. The information loss is geometric: the manifold $\mathbb{R}^L$ is partitioned into $M$ equivalence classes, and collapse corresponds to projecting the state onto one of these classes.

We now formulate this collapse more rigorously in geometric terms.

\begin{definition}
Let $X$ be a smooth manifold and let $\{\Omega_{k}\}$ be a measurable partition induced by a quantizer. The \emph{quantizer collapse map} is the canonical projection
\[
\pi_{Q} : X \to \Lambda,
\qquad
\pi_{Q}(x) = q_{k} \quad \text{if } x \in \Omega_{k}.
\]
\end{definition}

From the viewpoint of derived geometry \cite{toen2008homotopical, lurie2009higher}, collapse can be interpreted as a truncation functor. The quantizer forgets all homotopy information within each region $\Omega_{k}$ and retains only the connected-component label. Although the modulator’s state resides in $\mathbb{R}^{L}$, which is contractible, the quantizer regions produce a nontrivial quotient space with discrete topology.

The reconstruction map is the other half of the measurement process. After collapse, the loop filter updates the internal state by applying the vector field
\[
x[n+1] = A x[n] + B u[n] - B Q(x[n]).
\]
This vector field smooths the collapse by integrating the difference between the input and the symbolic output. The integrator serves as a conditional expectation operator that pulls the collapsed state back into a smooth manifold. The loop filter thereby provides a reconstruction map that restores coherence after measurement.

We now state the first central theorem of the chapter.

\begin{theorem}
The delta--sigma loop implements a deterministic classical measurement--reconstruction process. The quantizer $Q$ is a deterministic measurement in the sense of Busch--Lahti--Mittelstaedt, and the loop filter acts as a reconstruction operator equivalent to a conditional expectation with respect to the sigma-algebra generated by $Q$.
\end{theorem}

\begin{proof}
The quantizer defines a measurable partition of the state space into regions $\Omega_{k}$, each corresponding to a symbol $q_{k}$. The map $Q : X \to \Lambda$ is measurable and therefore constitutes a deterministic measurement. The sigma-algebra generated by $Q$, denoted $\sigma(Q)$, consists of unions of the regions $\Omega_{k}$. The loop filter update
\[
x[n+1] = A x[n] + B u[n] - B Q(x[n])
\]
is measurable with respect to $\sigma(Q)$ because it depends only on $x[n]$ and the symbol $Q(x[n])$. The integrator computes a smoothed version of the collapsed state by integrating the difference between $u[n]$ and $Q(x[n])$. This is equivalent to taking the conditional expectation of the uncollapsed update with respect to $\sigma(Q)$ since the collapsed state is the only available information. The reconstruction operator is therefore a conditional expectation, completing the measurement--reconstruction cycle.
\end{proof}

\noindent
This theorem places the delta--sigma modulator within the general measurement framework. The quantizer is a measurement, and the integrator is a reconstruction.

In the next section we deepen this connection by showing that the modulator realizes a classical analogue of the Lüders rule, where collapse is followed by a smoothing operation that preserves consistency between the state and the measurement outcome.

\section*{12.2 The Lüders–Rule Analogue for Delta--Sigma Modulation}

The connection between quantization and measurement becomes sharper when we examine the reconstruction step induced by the loop filter. In quantum measurement theory, the Lüders rule \cite{vonneumann1932foundations} replaces the pre-measurement density operator $\rho$ with the post-measurement state
\[
\rho \longmapsto \frac{P_{k} \rho P_{k}}{\mathrm{Tr}(P_{k} \rho)},
\]
where $P_{k}$ is the projection associated with the measurement outcome $k$. This update enforces consistency between the state and the observed outcome and performs a smoothing operation: coherence is preserved within the subspace onto which $P_{k}$ projects, and all incompatible amplitudes are removed.

In a classical system, the notion of collapse must be reinterpreted. Instead of projecting onto a subspace of a Hilbert space, collapse projects the state $x[n]$ onto the measurable partition associated with $Q$. The reconstruction step then restores coherence by integrating the discrepancy between the input and the collapsed symbolic state. This dual operation is the classical analogue of the Lüders rule.

We formalize this analogy in the following theorem.

\begin{theorem}
Let $Q$ be a deterministic quantizer and let $F$ denote the loop filter update
\[
x[n+1] = A x[n] + B u[n] - B Q(x[n]).
\]
Then the composite map
\[
x[n] \longmapsto Q(x[n]) \longmapsto x[n+1]
\]
is a classical analogue of the Lüders measurement update. It consists of a collapse onto the partition cell $\Omega_{k}$, followed by a reconstruction step that preserves coherence with the measurement outcome and minimizes the prediction error subject to the collapsed information.
\end{theorem}

\begin{proof}
The first step $x[n] \mapsto Q(x[n])$ is a collapse onto $\Omega_{k}$, in the sense that all points in $\Omega_{k}$ yield the same observable outcome. The second step applies the vector field $F$, which is affine in $x[n]$ and linear in $Q(x[n])$. This map ensures that $x[n+1]$ is consistent with the measurement outcome because the term $B Q(x[n])$ forces the new state to incorporate the information that the quantizer has collapsed. The reconstruction preserves coherence by smoothing the discontinuity introduced by the collapse through the integration term $A x[n] + B u[n]$, which is the discrete analogue of a conditional expectation with respect to the sigma-algebra $\sigma(Q)$. This yields a classical update rule that performs the two essential functions of the Lüders rule: collapsing onto the measurement-determined region and then restoring coherence under the constraints imposed by that collapse.
\end{proof}

This theorem establishes the first half of the measurement--reconstruction duality. The other half lies in the structure of the entropy field and the smoothing effect of the integrator, which we now analyze.

\section*{12.3 Entropy Collapse and Lamphrodynamic Smoothing}

The RSVP viewpoint interprets entropy not as a scalar quantity of uncertainty but as a dynamical field defined over the state space. Collapse under measurement injects entropy into the system because the projection $x[n] \mapsto Q(x[n])$ discards geometric information about the state. The integrator reduces entropy because it integrates a discrepancy between the input and the collapsed symbol, producing a more regular state trajectory.

Let $S(x)$ denote the entropy field defined in Chapter~7. Collapse induces an instantaneous increase in entropy:
\[
S(x[n]) < S\!\left(\Pi_{\Omega_{k}}(x[n])\right),
\]
where $\Pi_{\Omega_{k}}$ denotes projection onto the quantizer region. By contrast, the loop filter induces entropy reduction:
\[
S(x[n+1]) \leq S(x[n]),
\]
whenever the vector field is contractive in the sense of dissipative systems theory \cite{willems1972dissipative, khalil2002nonlinear}. This alternation between instantaneous entropy injection and gradual entropy reduction yields the characteristic error shaping of delta--sigma modulation: high-frequency entropy accumulates while low-frequency coherence is preserved.

We now formalize this behavior in a lamphrodynamic theorem.

\begin{theorem}
Let $S$ be an entropy field satisfying the lamphrodynamic inequality
\[
\frac{d}{dt}S(x(t)) \leq -\alpha \|v(x(t))\|^{2},
\]
for the vector field $v$ associated with the loop filter, where $\alpha > 0$. Then the delta--sigma modulator alternates between collapse-induced entropy increases and vector-field-induced entropy decreases. The long-term result is an entropy spectrum in which entropy is shaped away from the frequency band where $v$ possesses strong contraction.
\end{theorem}

\begin{proof}
Collapse under measurement produces a discontinuity in the state:
\[
x[n] \longmapsto Q(x[n]).
\]
Since $Q$ projects onto a discrete manifold, the entropy field $S$ increases under this mapping. Between measurements, the state evolves according to the vector field $v$, which satisfies the lamphrodynamic inequality. This inequality implies that entropy decreases monotonically during smooth evolution because high gradients of the vector field dissipate rapidly. Combining these observations over many samples yields an alternation between instantaneous entropy increases and smooth entropy decreases. The cumulative effect is spectral entropy shaping: entropy accumulates primarily in high-frequency components, while low-frequency components are suppressed. This behavior precisely matches the classical noise-shaping phenomenon described in \cite{schreier2005understanding} and admits a rigorous field-theoretic interpretation.
\end{proof}

\section*{12.4 Derived Geometry of Collapse and Reconstruction}

Collapse acts as a truncation functor in the sense of derived geometry \cite{lurie2009higher}. The continuous manifold $\mathbb{R}^{L}$ is replaced by a discrete set of partition labels. The integrator serves as a left adjoint to this truncation by embedding the discrete symbol back into the continuous manifold through the update equation.

\begin{definition}
Let $X$ be a smooth manifold and $\Lambda$ a discrete space. The quantizer $Q : X \to \Lambda$ induces a functor
\[
Q^{\sharp} : \mathbf{Sh}(X) \to \mathbf{Sh}(\Lambda),
\]
that pulls back sheaves along $Q$. The left adjoint $Q_{\sharp}$ reconstructs a sheaf on $X$ from a sheaf on $\Lambda$.
\end{definition}

In this language, collapse corresponds to $Q^{\sharp}$, while reconstruction corresponds to $Q_{\sharp}$. The delta--sigma loop thereby becomes a categorical adjunction:
\[
Q_{\sharp} \dashv Q^{\sharp}.
\]

We now prove the central categorical theorem of the chapter.

\begin{theorem}
The delta--sigma modulator implements the adjunction $Q_{\sharp} \dashv Q^{\sharp}$ at the level of state evolution. Collapse applies $Q^{\sharp}$ to the state, while the loop filter applies $Q_{\sharp}$. The composite update $Q_{\sharp} Q^{\sharp}$ is idempotent and defines a derived fiber product between the collapsed symbol and the reconstructed state.
\end{theorem}

\begin{proof}
Collapse maps $x[n] \in X$ to $Q(x[n]) \in \Lambda$, which is the action of $Q^{\sharp}$. The reconstruction update uses the symbol $Q(x[n])$ to compute $x[n+1]$ via the vector field $F$, which is the action of $Q_{\sharp}$. The composite map $x[n] \mapsto x[n+1]$ depends only on $Q(x[n])$. Therefore the composite $Q_{\sharp}Q^{\sharp}$ is idempotent:
\[
(Q_{\sharp}Q^{\sharp})^{2} = Q_{\sharp}Q^{\sharp}.
\]
This idempotent operator determines a derived fiber product between the discrete symbol and the continuous state because it reconstructs a canonical representative of the equivalence class determined by $Q$. The state update is therefore a derived geometric operation.
\end{proof}

This categorical structure unifies the measurement and reconstruction aspects of the delta--sigma modulator and connects it to higher structures in derived geometry.

\section*{12.5 Measurement--Reconstruction Duality: The Master Theorem}

We now state and prove the master theorem that unifies the concepts developed in this chapter.

\begin{theorem}[Measurement--Reconstruction Duality]
A delta--sigma modulator is equivalent to a classical measurement device followed by a reconstruction device. Quantization constitutes a deterministic measurement in the sense of Busch--Lahti--Mittelstaedt, while the loop filter implements a reconstruction operator that is both a conditional expectation and a left adjoint functor. Together they form an idempotent adjunction whose repeated application generates the coherent state trajectory characteristic of delta--sigma modulation.
\end{theorem}

\begin{proof}
The quantizer defines a deterministic measurement by collapsing the continuous state onto a discrete symbol. The loop filter update functions as a reconstruction because it uses the collapsed symbol to generate a new continuous state that is consistent with the measurement. The reconstruction is equivalent to conditional expectation because it minimizes the prediction error under the sigma-algebra generated by the quantizer. The categorical interpretation then shows that collapse and reconstruction form an adjunction, and idempotence follows because successive applications of collapse and reconstruction yield the same class of states. This adjunction generates the coherent, low-frequency structure of the state trajectory and shapes entropy into high-frequency components. Thus, the delta--sigma modulator is a measurement--reconstruction engine.
\end{proof}

The consequences of this theorem will be central in later chapters, where the modulator is embedded into the broader RSVP framework as an entropy routing device and a measurement primitive within lamphrodynamic field systems.

\section*{12.6 Vector-Valued Measurement Operators and Multi-Dimensional Collapse}

Up to this point, the analysis has focused primarily on scalar quantizers, where the measurement map
\[
Q : \mathbb{R}^{L} \to \Lambda
\]
projects the state onto a discrete set of scalar outputs. However, modern delta--sigma modulators frequently employ vector-valued quantizers, either because multiple quantizers operate in parallel, or because the modulator processes a multi-channel signal whose components are jointly quantized. Vector quantization is a deep subject in information theory \cite{gersho1991vector}, and its geometric structure greatly enriches the measurement--reconstruction duality.

Let the quantizer now be a map
\[
Q : X \subset \mathbb{R}^{L} \longrightarrow \Lambda \subset \mathbb{R}^{m},
\]
where $m \geq 1$ is the dimension of the output symbol. The measurement partitions the state space into regions
\[
\{\Omega_{\lambda}\}_{\lambda \in \Lambda},
\]
where each $\Omega_{\lambda} = Q^{-1}(\{\lambda\})$ is now a polyhedral or curved region in $\mathbb{R}^{L}$. Unlike the scalar case, these regions may have codimension greater than one, and their boundaries may intersect in highly nontrivial ways. The collapse $x \mapsto Q(x)$ therefore discards a richer set of geometric information.

The reconstruction step now depends on a vector-valued symbol $\lambda$:
\[
x[n+1] = A x[n] + B u[n] - B \lambda,
\]
where $\lambda = Q(x[n])$. The geometry of the reconstruction operator reflects the dimension and structure of $\Lambda$. If $\Lambda$ forms a lattice, then reconstruction resembles returning from a discrete grid to a continuous submanifold. If $\Lambda$ is irregular, the reconstruction must interpolate between nonuniformly spaced points, making the smoothing problem more delicate.

We formalize these ideas in the following theorem.

\begin{theorem}
Let $Q : X \to \Lambda$ be a vector-valued quantizer whose partition $\{\Omega_{\lambda}\}$ satisfies the regularity condition that each $\Omega_{\lambda}$ is bounded and has Lipschitz boundary. Then for any stable loop filter $(A,B)$, the delta--sigma reconstruction map
\[
F(x) = A x + B u - B Q(x)
\]
provides a Lipschitz-continuous smoothing operator on the quotient space $X/\!\sim_{Q}$, where $x \sim_{Q} x'$ if and only if $Q(x) = Q(x')$.
\end{theorem}

\begin{proof}
The collapse operator $Q$ induces a quotient topology on $X$, identifying points in the same region $\Omega_{\lambda}$. The reconstruction map $F$ depends continuously on the symbol $\lambda$ and affinely on the state $x$. Since $A$ is linear and $B$ is constant, $F$ is Lipschitz on each $\Omega_{\lambda}$ with Lipschitz constant $\|A\|$. Across boundaries, the change in $Q$ shifts $F$ by $B(\lambda - \lambda')$, which is bounded because $\Lambda$ is finite or discrete with bounded step size. The boundedness and Lipschitz nature of the partition boundaries ensure that these shifts do not disrupt continuity in the quotient topology. Therefore, $F$ defines a smoothing operator on $X/\!\sim_{Q}$.
\end{proof}

This theorem shows that vector-valued quantization yields a higher-dimensional version of the classical measurement--reconstruction duality, with the quotient space functioning as a multi-dimensional symbolic manifold.

\section*{12.7 Non-Commutative Quantizers and Measurement Geometry}

In many practical modulators, especially MIMO systems or modulators with multiple nested quantizers, the quantization operations may not commute. If $Q_{1}$ and $Q_{2}$ are two quantizers acting on overlapping subsets of the state space, their compositions may satisfy
\[
Q_{1} Q_{2} \neq Q_{2} Q_{1}.
\]
This non-commutativity introduces profound geometric consequences, analogous to the non-commutativity of observables in quantum mechanics.

Let
\[
Q_{1} : X \to \Lambda_{1}, \qquad Q_{2} : X \to \Lambda_{2},
\]
and define the composite measurement maps $Q_{12} = Q_{1} \circ Q_{2}$ and $Q_{21} = Q_{2} \circ Q_{1}$. If $Q_{12} \neq Q_{21}$, then the partitions they induce differ:
\[
\{\Omega_{12,\lambda}\} \neq \{\Omega_{21,\mu}\}.
\]
Thus the act of quantizing in different orders yields different quotients of the state space.

We capture this structure in the following lemma.

\begin{lemma}
Let $Q_{1}$ and $Q_{2}$ be two quantizers whose partitions do not refine one another. Then the iterated collapse maps $Q_{12}$ and $Q_{21}$ generate a non-commutative measurement algebra in which the sequence of quantizer applications determines the ultimate quotient of the state space.
\end{lemma}

\begin{proof}
If neither partition refines the other, there exist points $x$ and $x'$ such that $Q_{1}(x) = Q_{1}(x')$ but $Q_{2}(x) \neq Q_{2}(x')$. Applying $Q_{2}$ before $Q_{1}$ collapses such distinctions. Conversely, applying $Q_{1}$ first preserves them. Hence the compositions produce distinct results:
\[
Q_{1}(Q_{2}(x)) \neq Q_{2}(Q_{1}(x)).
\]
This implies that the two collapse maps act differently on equivalence classes of the induced quotient spaces. Therefore, the two quantizers generate a non-commutative algebra of state reductions.
\end{proof}

The implications for delta--sigma modulation are substantial, especially in multi-channel systems where different quantizers observe different combinations of the state vector. Non-commutative collapse can produce dynamical artifacts reminiscent of joint-measurement interference effects, although in a purely classical setting.

\section*{12.8 Cross-Channel Collapse and the Geometry of Coupled Modulators}

When multiple delta--sigma loops operate in parallel and are coupled through either shared state variables or shared quantizers, the measurement geometry acquires an additional layer of structure. Let $x^{(1)}[n]$ and $x^{(2)}[n]$ denote the states of two coupled modulators with a shared quantizer $Q$ acting on a linear combination of their states:
\[
y[n] = Q(\alpha x^{(1)}[n] + \beta x^{(2)}[n]).
\]
Collapse now affects both state vectors simultaneously, and reconstruction is performed independently by each modulator’s loop filter. This induces a correlated collapse process and a pair of correlated reconstruction maps.

If the coupling is symmetric, meaning $\alpha = \beta$, the system behaves like a multi-observer network performing a consensus operation under quantized communication. If the coupling is asymmetric, collapse can destabilize one modulator even if the other remains stable.

We formalize this in the following theorem.

\begin{theorem}
Let two delta--sigma modulators be coupled through a shared quantizer as above. If the loop filters of both modulators are stable in isolation, then the coupled system is stable provided that the coupling coefficients satisfy $|\alpha| = |\beta|$ and the loop filter matrices commute. If either condition is violated, collapse-induced divergence may occur.
\end{theorem}

\begin{proof}
When $|\alpha| = |\beta|$, the shared quantizer observes the arithmetic mean of the two states up to sign. The collapse therefore affects both states symmetrically. If the loop filter matrices commute, then the reconstruction maps preserve this symmetry. Under these conditions, collapse and reconstruction preserve the diagonal subspace $\{x^{(1)} = x^{(2)}\}$, and stability reduces to stability of the common mode. If either symmetry condition is violated, collapse affects the modulators differently, and their reconstruction maps no longer preserve a shared invariant subspace. This leads to the possibility that one state diverges while the other remains stable.
\end{proof}

This theorem shows that measurement geometry can induce or destroy stability in coupled delta--sigma networks.

\section*{12.9 Spectral Theory of Collapse: A Koopman Operator Approach}

Delta--sigma collapse may also be analyzed through the spectral theory of dynamical systems via the Koopman operator. The Koopman operator $\mathcal{K}$ associated with the map $x[n] \mapsto Q_{\sharp}Q^{\sharp}(x[n])$ acts on observables $f$ as
\[
(\mathcal{K} f)(x) = f(Q_{\sharp}Q^{\sharp}(x)).
\]
Because $Q_{\sharp}Q^{\sharp}$ is idempotent, the spectrum of $\mathcal{K}$ consists of eigenvalues $\{0,1\}$. The eigenspace associated with eigenvalue $1$ contains observables that are constant on each quantizer region, while the eigenspace associated with eigenvalue $0$ contains observables that vanish upon collapse.

This dichotomy reflects the fundamental decomposition of the state space into measurable and unmeasurable components. The integrator then acts as a smoothing operator on the measurable subspace, driving the system toward coherence.

We capture this insight in the following proposition.

\begin{proposition}
The Koopman operator of the collapse--reconstruction map decomposes the space of observables into invariant subspaces corresponding to measurable invariants and collapse-sensitive fluctuations. The delta--sigma loop shapes the spectral energy of observables so that fluctuations associated with eigenvalue $0$ accumulate at high frequencies, while invariants associated with eigenvalue $1$ remain coherent at low frequencies.
\end{proposition}

\begin{proof}
Idempotence of $Q_{\sharp} Q^{\sharp}$ implies that the Koopman operator is a projection. Projections have eigenvalues $0$ and $1$ only. The invariant subspace for eigenvalue $1$ contains functions constant on quantizer cells, while the eigenvalue $0$ subspace contains functions that detect differences within each cell. Collapse annihilates the latter, while reconstruction smooths the former. The integrator amplifies low-frequency components of the eigenvalue-$1$ observables, while the quantizer repeatedly suppresses eigenvalue-$0$ components, pushing their spectral energy into high-frequency components. This yields the characteristic spectral envelope of delta--sigma noise shaping.
\end{proof}

This spectral viewpoint provides a mathematically rigorous perspective on the phenomenon of quantization noise shaping and connects it to the operator-theoretic foundations of dynamical systems.

\section*{12.10 The Geometry of Reconstruction Manifolds}

The reconstruction map
\[
R(x) = A x + B u - B Q(x)
\]
defines a manifold structure on the state space. For each symbol $\lambda \in \Lambda$, define
\[
M_{\lambda} = \{ R(x) : Q(x) = \lambda \}.
\]
These sets form reconstruction manifolds, which encode how the modulator rebuilds the continuous state from symbolic information. Under certain regularity conditions, each $M_{\lambda}$ is diffeomorphic to the quantizer region $\Omega_{\lambda}$.

We state this precisely as a final theorem of this part.

\begin{theorem}
If $Q$ is locally constant on $\Omega_{\lambda}$ and $A$ is a diffeomorphism, then each reconstruction manifold $M_{\lambda}$ is diffeomorphic to $\Omega_{\lambda}$. Moreover, the collection of manifolds $\{M_{\lambda}\}$ forms a stratification of the state space that is invariant under the delta--sigma dynamics.
\end{theorem}

\begin{proof}
If $Q$ is constant on $\Omega_{\lambda}$, then for all $x \in \Omega_{\lambda}$, $R(x) = A x + B u - B \lambda$. Because $A$ is a diffeomorphism, the map $x \mapsto A x + B u - B \lambda$ is a diffeomorphism from $\Omega_{\lambda}$ to its image $M_{\lambda}$. The invariance of $\{M_{\lambda}\}$ under the delta--sigma update follows from the fact that $R$ maps $\Omega_{\lambda}$ into $M_{\lambda}$ and that the quantizer partitions drive all states back into some $\Omega_{\lambda'}$ at the next collapse, preserving the stratification.
\end{proof}

This geometric structure is critical for understanding multi-layer, multi-bit, and adaptive delta--sigma modulators developed in later chapters.

\section*{12.11 Stability Implications of Measurement-Based Reconstruction}

The measurement--reconstruction duality developed in earlier sections has profound consequences for stability. Classical stability theory for delta--sigma modulators relies on Lyapunov methods, dissipativity, and small-gain conditions \cite{khalil2002nonlinear, vidyasagar1993nonlinear}. The field-theoretic interpretation deepens this understanding by making explicit the competing effects of collapse and smoothing. Collapse always injects entropy, while smoothing always removes it. Stability arises when the smoothing effect dominates in the long-term average.

Formally, let $V : X \to \mathbb{R}_{\geq 0}$ be a candidate Lyapunov function and let $F$ denote the reconstruction operator. The collapse operator $C$ is defined by $C(x) = Q(x)$, and the combined update is
\[
x[n+1] = F(C(x[n])).
\]
We are interested in the behavior of $V(x[n])$ as $n \to \infty$. Collapse alone often increases $V$, while $F$ tends to decrease it.

This competition is encoded in the following inequality:
\[
V(F(C(x))) - V(x) = \underbrace{V(C(x)) - V(x)}_{\text{collapse-induced growth}} + \underbrace{V(F(C(x))) - V(C(x))}_{\text{reconstruction-induced decay}}.
\]
The second term is typically negative because $F$ performs smoothing. The first term is nonnegative because $C$ discards information. Stability requires that the second term eventually dominate the first.

We make this precise in the following theorem.

\begin{theorem}
Let $C$ be a quantizer and $F$ a reconstruction map defined by a stable loop filter. Suppose there exist constants $\alpha > 0$ and $\beta \geq 0$ such that for all $x$,
\[
V(F(C(x))) - V(C(x)) \leq -\alpha \|C(x)\|^{2},
\]
and suppose collapse satisfies the growth bound
\[
V(C(x)) - V(x) \leq \beta.
\]
Then the delta--sigma modulator is globally bounded-input bounded-state stable, and the state satisfies
\[
\limsup_{n \to \infty} V(x[n]) \leq \beta / \alpha.
\]
\end{theorem}

\begin{proof}
We write
\[
V(x[n+1]) - V(x[n]) = V(F(C(x[n]))) - V(C(x[n])) + V(C(x[n])) - V(x[n]).
\]
By the assumptions,
\[
V(x[n+1]) - V(x[n]) \leq -\alpha \|C(x[n])\|^{2} + \beta.
\]
Summing over $n$ yields
\[
\sum_{k=0}^{N} \alpha \|C(x[k])\|^{2} \leq V(x[0]) - V(x[N+1]) + N \beta.
\]
Since $V$ is non-negative, we obtain
\[
\frac{1}{N+1} \sum_{k=0}^{N} \|C(x[k])\|^{2} \leq \frac{V(x[0])}{\alpha(N+1)} + \frac{\beta}{\alpha}.
\]
Taking the limit as $N \to \infty$ yields
\[
\limsup_{N \to \infty} \frac{1}{N+1} \sum_{k=0}^{N} \|C(x[k])\|^{2} \leq \frac{\beta}{\alpha}.
\]
This implies boundedness of the squared norm of the quantizer outputs and therefore boundedness of the state because the reconstruction map $F$ is affine and $A$ is stable. Thus $x[n]$ is bounded for all $n$, establishing bounded-input bounded-state stability. The upper bound on $\limsup V(x[n])$ follows immediately.
\end{proof}

This theorem generalizes classical stability results by incorporating the measurement--reconstruction viewpoint and is robust to nonlinear, piecewise-affine, or discontinuous quantizers.

\section*{12.12 Measurement Geometry and Catastrophic Instability}

Measurement-induced instability is a well-known phenomenon in higher-order delta--sigma modulators. Traditionally it is analyzed through root-locus or describing-function techniques \cite{schreier2005understanding}, but the geometric interpretation developed here provides a more unified view.

Collapse partitions the state space into regions $\Omega_{\lambda}$, and the reconstruction map $F$ maps each region into a reconstruction manifold $M_{\lambda}$. Stability requires that trajectories move between these manifolds in a coherent fashion. Catastrophic instability arises when the reconstruction map sends a state deep into a region from which collapse drives it into a reconstruction manifold that induces further divergence rather than contraction.

We formalize this phenomenon.

\begin{theorem}
Let the reconstruction manifolds $\{M_{\lambda}\}$ be invariant under $F$. Catastrophic instability occurs if and only if there exist $\lambda$ and $\lambda'$ such that $F(M_{\lambda}) \subset \Omega_{\lambda'}$ and the induced map $F : M_{\lambda} \to M_{\lambda'}$ has a dominant expanding direction whose contraction by collapse is insufficient to offset its growth.
\end{theorem}

\begin{proof}
If $F(M_{\lambda}) \subset \Omega_{\lambda'}$, then collapse sends the image of $M_{\lambda}$ into $\Omega_{\lambda'}$ and reconstruction sends it into $M_{\lambda'}$. The combined map $T = F \circ C$ maps $M_{\lambda}$ into $M_{\lambda'}$. If $DF(x)$ has an eigenvalue greater than $1$ in magnitude along some direction and if $C$ does not contract this direction because points in this direction lie within the same quantizer cell, then $DT(x)$ inherits the expanding direction. Repeated application of $T$ therefore yields exponential divergence. Conversely, if all expanding directions are contracted by collapse, no instability can occur. Thus instability arises precisely when an expanding direction is invisible to the quantizer.
\end{proof}

This theorem explains why higher-order modulators with coarse quantizers are prone to instability: the quantizer’s partition is too coarse to suppress expanding directions in the loop dynamics.

\section*{12.13 Measurement Sensitivity and Quantizer Geometry}

The shape of the quantizer regions $\Omega_{\lambda}$ influences the sensitivity of the modulator to small perturbations. If the quantizer boundaries are nearly tangent to a direction of high sensitivity in the loop filter, then small perturbations in that direction may cause large differences in collapse outcomes.

Let $n_{\partial \Omega}$ be the normal vector to the quantizer boundary. If $v(x)$ is the vector field induced by the loop filter, collapse sensitivity is determined by the angle between $v(x)$ and $n_{\partial \Omega}$:
\[
\theta(x) = \arccos \left( \frac{\langle v(x), n_{\partial \Omega} \rangle}{\|v(x)\| \|n_{\partial \Omega}\|} \right).
\]
If $\theta(x) \approx \pi/2$, then the trajectory grazes the boundary, and small variations in $x$ produce abrupt quantizer transitions. This grazing phenomenon is responsible for idle tones and periodic patterns in many modulators.

We state the following lemma.

\begin{lemma}
Collapse sensitivity is maximized when the loop filter vector field is nearly tangent to the quantizer boundary. In this case, small perturbations in the tangential direction are preserved under reconstruction but produce large variations in collapse outcome.
\end{lemma}

\begin{proof}
If $v(x)$ is tangent to $\partial \Omega_{\lambda}$, then the trajectory remains near the boundary for multiple iterations. Perturbations in the tangent direction produce large variations in crossing times. Collapse is determined by the sign of the normal component of the motion, which becomes extremely sensitive to perturbations. Since reconstruction preserves tangential perturbations, measurement outcomes diverge rapidly, maximizing collapse sensitivity.
\end{proof}

Understanding collapse sensitivity is essential for designing quantizers whose geometry is compatible with stable loop dynamics.

\section*{12.14 Entropy Allocation and Measurement Potential}

Every quantizer defines an entropy potential
\[
H_{Q}(x) = \log |\Omega_{Q(x)}|,
\]
where $|\Omega_{Q(x)}|$ denotes the Lebesgue measure of the quantizer region. Collapse increases $H_{Q}$ because regions with larger measure absorb more initial states. The loop filter tends to decrease $H_{Q}$ because it contracts states into narrower regions.

The entropy allocation reveals deep properties of modulator efficiency. The optimal quantizer, from an entropy perspective, is one whose regions have equal measure when weighted by the invariant distribution of the loop dynamics. This is analogous to the Shannon condition for optimal scalar quantization \cite{gersho1991vector}.

We state a final proposition.

\begin{proposition}
A delta--sigma quantizer minimizes long-term entropy production if and only if its partition regions have equal measure with respect to the invariant measure of the collapse--reconstruction map.
\end{proposition}

\begin{proof}
Entropy production per collapse is $H_{Q}(x[n])$. Averaging over the invariant measure $\mu$ of the map yields expected entropy production
\[
\mathbb{E}_{\mu}[H_{Q}] = \sum_{\lambda} \mu(\Omega_{\lambda}) \log |\Omega_{\lambda}|.
\]
Using Jensen's inequality and the convexity of the logarithm, this quantity is minimized when all $|\Omega_{\lambda}|$ are equal under the weighting $\mu(\Omega_{\lambda})$, yielding the stated condition.
\end{proof}

This proposition provides a geometric criterion for designing optimal quantizers in delta--sigma systems, generalizing classical results to the collapse--reconstruction framework.

\section*{12.15 The PDE Limit of Collapse--Reconstruction Dynamics}

The alternating sequence of collapse and reconstruction can be approximated in the high-OSR limit by a partial differential equation whose structure reflects both the smoothing of the loop filter and the singular impulses introduced by the quantizer. This limit is central to the unification of delta--sigma modulation with lamphrodynamic flow theory, and it yields a rigorous continuum interpretation of noise shaping.

Let $\Delta t$ be the sampling period of the digital or switched-capacitor implementation, and define the scaled state trajectory $x(t)$ by $x[n] = x(n \Delta t)$. In the limit $\Delta t \to 0$, the discrete map
\[
x[n+1] = A x[n] + B u[n] - B Q(x[n])
\]
approaches a differential equation whose right-hand side is piecewise smooth. Write the reconstruction operator as
\[
R(x) = A x + B u - B Q(x).
\]
Expanding in $\Delta t$,
\[
x(t+\Delta t) - x(t) = \Delta t \, \frac{dx}{dt} + O(\Delta t^{2}),
\]
and equating both sides,
\[
\frac{dx}{dt} = A x + B u - B Q(x) + O(\Delta t).
\]
Taking the limit yields the piecewise-smooth ODE
\[
\frac{dx}{dt} = A x + B u - B Q(x).
\]

However, this expression conceals the singular nature of collapse. To capture collapse explicitly, write the quantizer as
\[
Q(x) = x - e(x),
\]
where $e(x)$ takes values in a discrete set. Substituting into the ODE,
\[
\frac{dx}{dt} = A x + B u - B x + B e(x)
\]
or
\[
\frac{dx}{dt} = (A - B C)x + B u + B e(x),
\]
where $C$ is the identity in the scalar case. The term $B e(x)$ is a distribution-valued function in the continuum limit.

We formalize the PDE limit as follows.

\begin{theorem}
Let $x[n]$ evolve under the delta--sigma update with sampling period $\Delta t$. In the asymptotic limit $\Delta t \to 0$ with input held fixed, the state trajectory converges in the sense of distributions to a solution of the PDE
\[
\frac{\partial x}{\partial t} = v(x,u) + B \eta(x,t),
\]
where $v(x,u) = A x + B u - B C x$ is the lamphrodynamic vector field and $\eta(x,t)$ is a distribution-valued entropy injection term supported on the quantizer boundaries.
\end{theorem}

\begin{proof}
Writing the update as
\[
x[n+1] - x[n] = \Delta t \, v(x[n],u[n]) + \Delta t \, B e[x[n]],
\]
and dividing by $\Delta t$ yields
\[
\frac{x[n+1] - x[n]}{\Delta t} = v(x[n]) + B e[x[n]].
\]
As $\Delta t \to 0$, the left-hand side converges in the distributional sense to $\partial x / \partial t$. The term $v(x)$ is smooth. The quantizer error $e[x[n]]$ converges to a distribution supported at the times where the state crosses quantizer thresholds, which correspond to curves in state-time space. Denoting this distribution by $\eta(x,t)$ gives the stated PDE.
\end{proof}

This PDE explicitly encodes the alternation of smoothing via $v(x)$ and impulse-like entropy insertion via $\eta(x,t)$.

\section*{12.16 Weak Solutions and Entropy Jumps}

The delta--sigma PDE cannot in general be interpreted classically because the entropy injection term is measure-valued. The correct notion is that of a weak solution. Let $\varphi(x,t)$ be a smooth test function with compact support. The weak form of the PDE is
\[
\int_{0}^{T} \int_{\mathbb{R}^{L}} x \, \partial_{t} \varphi + v(x,u) \cdot \nabla \varphi \, dx \, dt
=
\int_{0}^{T} \int_{\mathbb{R}^{L}} B \eta(x,t) \, \varphi(x,t) \, dx \, dt.
\]

The right-hand side integrates contributions from all threshold crossings. These contributions are precisely the entropy jumps induced by collapse.

We make this precise in the following lemma.

\begin{lemma}
Let $x(t)$ be a trajectory of the delta--sigma PDE. At any time $t_{0}$ where $x(t_{0}^{-})$ and $x(t_{0}^{+})$ lie in different quantizer cells, the entropy field satisfies
\[
S(x(t_{0}^{+})) - S(x(t_{0}^{-})) = \Delta S_{Q},
\]
where $\Delta S_{Q}$ is a constant determined by the quantizer geometry. These jumps contribute Dirac delta components to the distribution $\eta(x,t)$.
\end{lemma}

\begin{proof}
Collapse replaces $x(t_{0}^{-})$ by $Q(x(t_{0}^{-})) = x(t_{0}^{+})$. The entropy field $S$ increases by the entropy associated with discarding all geometric information in the pre-collapse equivalence class. In the continuum limit this instantaneous change becomes a distributional jump. The induced change in the PDE is therefore a Dirac delta supported at $t=t_{0}$.
\end{proof}

This weak formulation is central to understanding how quantization noise manifests as high-frequency entropy in the reconstructed signal.

\section*{12.17 The Continuum NTF as a Pseudodifferential Operator}

In the discrete-time theory, the noise transfer function
\[
\mathrm{NTF}(z) = 1 - C (zI - A)^{-1} B
\]
shapes quantization noise spectrally. In the continuum limit, the operator $(zI - A)^{-1}$ becomes the resolvent of the differential operator $\partial_{t} - A$. This transition leads to a pseudodifferential operator representation of noise shaping.

Let $L = \partial_{t} - A$ denote the lamphrodynamic differential operator. Then the continuum NTF is
\[
\mathcal{N} = B^{\dagger} L^{-1} C^{\dagger},
\]
where $L^{-1}$ is the Green’s operator for $L$.

We state the main PDE noise-shaping theorem.

\begin{theorem}
In the continuum limit, the delta--sigma noise-shaping operator is the pseudodifferential operator $\mathcal{N} = B^{\dagger} L^{-1} C^{\dagger}$, and the quantization noise $e(t)$ becomes a distribution supported on quantizer boundaries. The modulator output satisfies
\[
y(t) = \mathcal{S} u(t) + \mathcal{N} e(t),
\]
where $\mathcal{S}$ is the signal transfer operator.
\end{theorem}

\begin{proof}
Applying the Fourier transform to the discrete-time expression yields
\[
Y(\omega) = \mathrm{STF}(\omega)U(\omega) + \mathrm{NTF}(\omega)E(\omega).
\]
In the limit $\Delta t \to 0$, the variable $z$ approaches $e^{i\omega \Delta t}$, and
\[
z^{-1} \approx 1 - i\omega \Delta t.
\]
The matrix $(zI - A)^{-1}$ converges to $(i\omega I - A)^{-1}$, which is the Fourier-domain representation of $L^{-1}$. The remaining terms converge similarly. Therefore the limiting operator is the pseudodifferential operator stated above.
\end{proof}

This result unites the discrete frequency-domain theory with continuous field dynamics and establishes noise shaping as a smoothing of the lamphrodynamic PDE.

\section*{12.18 Entropy Flux Conservation in the Continuum Limit}

The flow of entropy in the PDE satisfies a conservation law away from the quantizer boundaries. Differentiating $S(x(t))$ along a smooth segment gives
\[
\frac{d}{dt} S(x(t)) = \langle \nabla S, v(x(t)) \rangle.
\]
If $v$ is lamphrodynamically contractive, then $dS/dt \leq 0$, implying that entropy decays during smooth evolution. At quantizer boundaries, entropy increases by $\Delta S_{Q}$. Summing over all boundaries crossed in $[0,T]$ yields
\[
S(x(T)) - S(x(0)) = -\int_{0}^{T} \langle \nabla S, v(x(t)) \rangle dt + \sum_{t_{k} \in [0,T]} \Delta S_{Q}.
\]

This balance law is the continuum analogue of the discrete smoothing-collapse alternation and demonstrates that the delta--sigma modulator enforces a global entropy routing principle.

\section*{12.19 The Lamphrodynamic Limit Cycle Theorem}

The PDE admits limit cycles corresponding to idle tones in discrete delta--sigma modulators. We now prove a lamphrodynamic criterion for the existence of such cycles.

\begin{theorem}
Let $x(t)$ evolve under the delta--sigma PDE. A limit cycle exists if and only if the entropy flux injected by quantizer crossings over one period exactly balances the entropy dissipated by lamphrodynamic smoothing:
\[
\oint_{\gamma} \langle \nabla S, v(x) \rangle dt + \sum_{k \in \gamma} \Delta S_{Q} = 0.
\]
\end{theorem}

\begin{proof}
A periodic solution satisfies $S(x(T)) = S(x(0))$. Writing the entropy balance over one period gives the identity above. If the sum is nonzero, entropy either accumulates or dissipates over successive cycles, contradicting periodicity. If the balance holds exactly, an entropy-neutral cycle exists, corresponding to a limit cycle in the discrete system.
\end{proof}

Idle tones are therefore entropy-balanced periodic orbits of the lamphrodynamic PDE.

\section*{12.20 Continuum Interpretation of Spectral Leakage}

Spectral leakage in delta--sigma modulation arises from the imperfect annihilation of low-frequency components of $e(t)$ by the pseudodifferential operator $\mathcal{N}$. In the PDE setting, this leakage corresponds to non-orthogonality between $\eta(x,t)$ and the nullspace of $L^{-1}$.

Let $\phi$ be a solution of $L \phi = 0$. The leakage energy is
\[
\int_{0}^{T} \phi(t) \, \mathcal{N} \eta(t) \, dt.
\]

We conclude this section with a proposition linking leakage to PDE geometry.

\begin{proposition}
Spectral leakage is proportional to the projection of the entropy injection distribution $\eta$ onto the approximate nullspace of the lamphrodynamic differential operator $L$. Minimizing leakage therefore requires aligning quantizer thresholds with level sets of the slow modes of $L$.
\end{proposition}

\begin{proof}
The operator $L^{-1}$ amplifies components of $\eta$ aligned with its nullspace. The nullspace corresponds to the slowest lamphrodynamic modes. Leakage appears precisely when $\eta$ has nonzero projection onto these modes. Aligning quantizer thresholds with level sets of slow modes minimizes this projection and hence minimizes leakage.
\end{proof}

This last result provides a continuous geometric interpretation of spectral leakage and guides the design of quantizers with improved spectral purity.

\section*{12.21 Variational Formulation of Collapse--Reconstruction Dynamics}

The lamphrodynamic interpretation of delta--sigma modulation becomes complete only when the dynamics are embedded in a variational framework. The alternation of smoothing and collapse corresponds to a system whose evolution minimizes a composite action functional consisting of a smooth Lagrangian contribution and a sequence of impulsive collapse terms. In this section we derive the continuum action functional, the Euler--Lagrange equations governing smooth flow, and the jump conditions induced by quantizer collapse.

Let the state be $x(t) \in \mathbb{R}^{L}$ evolving under a piecewise-smooth vector field $v(x,u)$ with entropy injection given by $\eta(x,t)$. We define the action functional
\[
\mathcal{A}[x] = \int_{0}^{T} \mathcal{L}(x(t),\dot{x}(t)) \, dt + \sum_{t_k \in \mathcal{C}} \mathcal{J}(x(t_k^-),x(t_k^+)),
\]
where $\mathcal{C}$ is the set of collapse times and $\mathcal{J}$ is the collapse penalty functional.

A natural choice for the Lagrangian is
\[
\mathcal{L}(x,\dot{x}) = \frac{1}{2} \|\dot{x} - v(x,u)\|_{M}^{2},
\]
where $M$ is a positive-definite inner-product matrix. This Lagrangian penalizes deviation from the lamphrodynamic vector field; its minimization enforces adherence to the vector field except at collapse events. The collapse penalty functional is defined by
\[
\mathcal{J}(x^{-},x^{+}) = S(x^{+}) - S(x^{-}),
\]
the entropy jump induced by quantizer collapse.

\begin{theorem}
Let $x(t)$ be a minimizer of the action $\mathcal{A}$. Then $x(t)$ satisfies the Euler--Lagrange equation
\[
M (\ddot{x} - Dv(x)\dot{x}) + M Dv(x)^{\top}( \dot{x} - v(x)) = 0
\]
on intervals between collapse events, and satisfies the collapse jump condition
\[
M (\dot{x}(t_k^+) - \dot{x}(t_k^-))
= \nabla_{x^{+}} \mathcal{J}(x(t_k^-),x(t_k^+)).
\]
\end{theorem}

\begin{proof}
The variational derivative of the integral term yields the Euler--Lagrange equation
\[
\frac{d}{dt}\left( M(\dot{x} - v(x)) \right)
- M Dv(x)(\dot{x} - v(x)) = 0.
\]
Expanding the derivative gives the stated form. The contributions from collapse times come from the jump terms. Variation of $\mathcal{J}$ with respect to $x^{+}$ yields the stated momentum-jump condition. The collapse times are fixed by the quantizer thresholds, so no time-variation terms occur.
\end{proof}

This establishes that the smooth segments of a delta--sigma trajectory follow geodesics of the lamphrodynamic metric modulated by the vector field $v(x)$.

\section*{12.22 Geometric Interpretation: Collapse as a Discrete Connection}

The Lamphrodynamic action induces a geometric structure on the state manifold. Between collapse events, the dynamics follow an affine connection with Christoffel symbols derived from the vector field Jacobian $Dv(x)$. Collapse induces an instantaneous change in the conjugate momentum $\pi = M (\dot{x} - v(x))$.

We now formalize this structure.

\begin{lemma}
Let $\pi(t) = M(\dot{x} - v(x))$ be the lamphrodynamic momentum. At collapse times $t_k$, the momentum satisfies
\[
\pi(t_k^+) = \pi(t_k^-) + \nabla S(x(t_k^+)).
\]
\end{lemma}

\begin{proof}
This follows immediately from the variational jump condition with $M$ as the mass matrix and collapse penalty $\mathcal{J}(x^-,x^+) = S(x^+) - S(x^-)$.
\end{proof}

Thus collapse acts as a discrete connection rule transforming the momentum by a gradient vector orthogonal to the quantizer manifold. From a differential geometric standpoint, this is analogous to a discrete holonomy operator acting on the cotangent bundle.

\section*{12.23 The Collapse Hypersurface and Its Normal Geometry}

Let $\Sigma_k = \{ x : Cx = q_k \}$ denote the quantizer hypersurface associated with output level $q_k$. The system undergoes collapse when the state crosses $\Sigma_k$. The normal vector is given by
\[
n_k = \frac{C^\top}{\|C\|}.
\]

At a collapse event the state is projected onto $\Sigma_k$:
\[
x(t_k^+) = Q(x(t_k^-)).
\]

We define the collapse map $\Pi_k : \mathbb{R}^{L} \to \Sigma_k$ by
\[
\Pi_k(x) = x - (Cx - q_k)\frac{C^\top}{\|C\|^{2}}.
\]

\begin{theorem}
The collapse map $\Pi_k$ is the metric projection onto $\Sigma_k$ with respect to the Euclidean metric induced by $M$ if and only if $MC^\top = \lambda C^\top$ for some scalar $\lambda$.
\end{theorem}

\begin{proof}
$\Pi_k$ is a metric projection onto the hyperplane $Cx = q_k$ precisely when $C^\top$ is an $M$-orthogonal normal to the hyperplane. This requires $M C^\top$ to be a scalar multiple of $C^\top$, establishing the condition.
\end{proof}

This characterization is important for understanding how quantizer design influences the geometry of collapse.

\section*{12.24 Hamiltonian Formulation and Energy Exchange at Collapse}

Define the lamphrodynamic Hamiltonian
\[
\mathcal{H}(x,\pi) = \frac{1}{2} \pi^\top M^{-1} \pi + \langle \pi, v(x) \rangle.
\]

Hamilton’s equations give
\[
\dot{x} = M^{-1} \pi + v(x), \qquad \dot{\pi} = - Dv(x)^\top \pi.
\]

Collapse induces an instantaneous change
\[
\pi(t_k^+) = \pi(t_k^-) + \nabla S(x(t_k^+)),
\]
which in turn produces a discrete energy jump
\[
\Delta \mathcal{H}(t_k) =
\langle \pi(t_k^-), \nabla S(x(t_k^+)) \rangle
+ \frac{1}{2} \| M^{-1/2}\nabla S(x(t_k^+)) \|^{2}.
\]

This energy jump is the Hamiltonian representation of quantization noise. The stochastic interpretation follows immediately.

\section*{12.25 The Entropy Routing Principle: A Variational Characterization}

The alternation of smoothing and collapse implements an entropy routing mechanism: entropy introduced at collapse is directed into high-frequency modes where it has minimal effect on reconstruction fidelity. We now derive a variational theorem capturing this principle.

\begin{theorem}[Entropy Routing Principle]
Let $x(t)$ minimize the lamphrodynamic action with collapse. Then the induced frequency distribution of quantization noise minimizes the functional
\[
\mathcal{E}[\eta] = 
\int_{0}^{T} 
\| L^{-1}\eta(t) \|_{W}^{2} \, dt
\]
subject to the constraint that $\eta$ consists of impulses supported on quantizer boundaries. Here $W$ is a positive-definite weighting matrix.
\end{theorem}

\begin{proof}
The operator $L^{-1}$ determines the effect of entropy injection on the trajectory. The action integral penalizes squared deviation from $v(x)$, which after applying the Green’s operator becomes a quadratic functional of $\eta$. The constraints enforce that $\eta$ is supported on threshold crossings. Minimization of the quadratic functional under these constraints yields the stated result.
\end{proof}

Thus the modulator performs optimal entropy routing given the constraint that collapse must occur along quantizer hypersurfaces.

\section*{12.26 Collapse as a Derived Geometric Truncation}

In derived geometry, quantizer collapse corresponds to a truncation functor. The state is a point in a derived fiber product
\[
X = \mathbb{R}^{L} \times_{\mathbb{R}} \Lambda,
\]
where $\Lambda$ is the discrete quantizer output set. Collapse corresponds to applying a truncation
\[
\tau_{\le 0} : X \to \Lambda,
\]
followed by a derived embedding $\iota : \Lambda \to \mathbb{R}^{L}$.

Thus the collapse--reconstruction cycle is the composition
\[
x \mapsto \tau_{\le 0}(x) \mapsto \iota\big(\tau_{\le 0}(x)\big).
\]

This derived interpretation explains why collapse introduces singularities into the lamphrodynamic PDE and why the entropy injection term is distribution-valued.

\section*{12.27 Summary of the Variational and PDE Framework}

The results of this chapter establish the following structure.  
First, delta--sigma modulation admits a rigorous PDE limit capturing its smoothing and collapse dynamics.  
Second, this PDE arises from a variational principle whose Euler--Lagrange equations govern lamphrodynamic segments and whose jump conditions describe collapse.  
Third, the entropy dynamics obey an exact conservation law with discrete jumps determined by quantizer geometry.  
Fourth, the noise-shaping operator in the continuum is a pseudodifferential Green’s operator associated with the lamphrodynamic differential operator.  
Finally, the entire system has a natural interpretation in derived geometry as alternating truncation and embedding on a derived fiber product.

This variational, geometric, and derived framework provides the foundational structure for all further analysis of delta--sigma modulators in measurement, control, and cognitive field theory.
